{%
\def\insertshortauthor{吕昀键}
\def\insertshorttitle{07.20\textasciitilde07.26 工作进展}
\def\insertshortdate{2026.07.27}
\begin{frame}[t]{低精度实验: 长度外推优势保持}
  \centering
  \vspace{0cm}
  \begin{minipage}{0.94\textwidth}
  \begin{columns}[T,onlytextwidth]
    \begin{column}{0.71\linewidth}
      \vskip0.08cm
      \makebox[\linewidth][c]{%
        \includegraphics[
          width=1.17\linewidth,
          height=5.15cm,
          keepaspectratio
        ]{20260720-0726/p19/matching-precision-last.png}%
      }
    \end{column}

    \begin{column}{0.26\linewidth}
      \thinsectionbar{完整覆盖}
      \vskip0.04cm
      \begin{beamercolorbox}[sep=3pt,rounded=true,wd=\linewidth]{block body}
        \scriptsize
        2080 Ti: FP32/FP16.\par
        RTX 3090: FP32/FP16/BF16.\par
        \textcolor{scured}{\bfseries 30/30} 个正式 4 epoch run.
      \end{beamercolorbox}

      \vskip0.06cm
      \thinsectionbar{质量与机制}
      \vskip0.04cm
      \begin{beamercolorbox}[sep=3pt,rounded=true,wd=\linewidth]{block body}
        \scriptsize
        外推切片: Flash \textcolor{scured}{\bfseries 60/60} 配对领先.\par
        训练端点: 仅 \textbf{3/15} 配对领先.\par
        \vskip0.05cm
        固定 checkpoint 改 eval dtype:\par
        max accuracy span $\leq \textbf{0.002328}$.
      \end{beamercolorbox}

      \vskip0.06cm
      \thinsectionbar{RTX 3090 Flash BF16}
      \vskip0.04cm
      \begin{beamercolorbox}[sep=3pt,rounded=true,wd=\linewidth]{block body}
        \scriptsize
        Wall time: \textcolor{scured}{\bfseries 0.862$\times$} FP32.\par
        Peak allocated: \textcolor{scured}{\bfseries 0.819$\times$} FP32.\par
        \vskip0.05cm
        主要差异来自低精度\textbf{训练轨迹}, 不是评估 dtype.
      \end{beamercolorbox}
    \end{column}
  \end{columns}

  \vskip0cm
  \begin{beamercolorbox}[sep=3pt,rounded=true,wd=\linewidth]{block body}
    \centering\small
    \textcolor{scured}{\bfseries 结论: FP16/BF16 未破坏 Flash 的长度外推优势; RTX 3090 BF16 兼具质量与效率依据.}
  \end{beamercolorbox}
  \end{minipage}
\end{frame}
}
